\documentclass{jsarticle}
\usepackage[dvipdfmx,final]{graphicx}
\usepackage{amsmath}
\usepackage{amssymb}
\begin{document}
\title{Farmat have with three number concluded from \\
gamma function of global partial integral manifold}
\author{Masaaki Yamaguchi}
\date{}
\maketitle

  $$ x^n + y^n \le z^n $$
  $$ \int (x^n + y^n)\mathrm{dvol} \le \int z^n dz_m $$
  $$ e^{x \log x} + e^{y \log y} \le e^{z \log z} $$
  $$ x \log x = u, y \log y = v, z \log z=w $$
  $$ u + y \ge z $$
  $$ n + n \ge n $$
  $$ n < 0, \int \Gamma(\gamma)^{'}dx_m = e^{- x \log x} $$
  ここで、$ n > 4 $だと、$ x^{u} \ge 0, y^{v} \ge, z^{w} \ge $となり、
  虚数をべき乗にとり、$ n < 3 $だと、$ e^{x \log x} + e^{y \log y} $
  が、$ x^x + y^y \le z^z $となり、$ (x + y)(x - y) = z $
  と、最小値に、メビウス空間となり、補空間に種数３をとり、ガロワ群となり、
  $ e^{x \log x} + e^{y \log y} \to e^{(x\log x)+(y \log y)} = e^{z \log z} $
  $$ e^{\log x^x + \log y^y} = e^{\log (x^x\cdot y^y)} \le e^{\log z^z} $$ 
  $$ x^x \cdot y^y = z^z $$
  整数の性質より、(1,1),(2,2),(3,3)であり、それ以外の整数は、虚数解をもつ
  $ e^{-f} $と　
  $$ e^{4+4}\ge 0 $$
  となり、$ z > 4 $だと、$x^x \cdot y^y \neq 0 $
  でなくなり、且つ、虚数が解に持たられる。
  また、べき乗に負が現れなくなる。
  $$ \int \Gamma(\gamma)^{'}dx_m = \Gamma^{\gamma} $$
  より、
  $$ x^{x^{{3}^{'}}} = x^{x^2} $$
  と言えて、
  $$ x=y=z \le 3 $$
  となる。
  もし、$$ z=4, z^{z^{{4}^{'}}} $$
  だと、
  $$ x + y = 4 $$
  となり、
  $ e^{-f} $と、負にならずに、$ x^x = e^{x \log x} $と、ガンマ関数による
  大域的部分積分多様体の解にならない。
  ガンマ関数による大域的部分積分多様体の解は、
  $$ \int \Gamma(\gamma)^{'}dx_m = e^{-x \log x} $$
  である。

\end{document}
